% Capítulo 6: Conclusiones y Trabajo Futuro

\section{Conclusiones}

Este trabajo ha evaluado experimentalmente Qsimov, una librería Python para re-entrenamiento
incremental sin olvido catastrófico en redes neuronales con activación ReLU. La contribución central
de este TFG es la caracterización experimental del comportamiento de Qsimov sobre MNIST, CIFAR-10
e ImageNet (100 clases), analizando las condiciones bajo las cuales su mecanismo de acumulación de
ecuaciones lineales elimina el olvido catastrófico.

Los resultados confirman que:

\begin{itemize}
  \item El sistema lineal de Qsimov supera a todos los métodos comparados en el escenario de
        aprendizaje continual sobre ImageNet, con un 37.0\,\% de precisión global frente al 23.0\,\%
        del re-entrenamiento acumulado y el 20.8\,\% del fine-tuning estándar.
  \item El tiempo de actualización del sistema lineal permanece prácticamente constante
        ($\approx$5.8\,s por batch en el experimento de streaming de ImageNet),
        independientemente del número de clases ya aprendidas.
  \item La robustez al learning rate del módulo de gradiente es notable en problemas de baja
        complejidad, eliminando la necesidad de ajuste fino de este hiperparámetro.
\end{itemize}

El principal factor limitante identificado es el número de caminos generados por el PathSelector,
que depende directamente del punto de corte elegido y de la densidad de la red base. En redes con
pocas neuronas antes de la capa de salida (CIFAR-10 con 33 caminos), el espacio representacional
puede no ser suficiente para aprender nuevas clases sin comprometer la capacidad en las antiguas.

Los resultados obtenidos son representativos de redes feed-forward con activación ReLU que incluyen
capas densas en su parte final, que es la arquitectura para la que Qsimov está diseñado. Cabe
esperar un comportamiento análogo en otras redes que cumplan estas condiciones, dado que el
mecanismo de acumulación de ecuaciones es agnóstico al dominio del problema. La extrapolación a
arquitecturas con conexiones residuales (ResNet, EfficientNet) o a modelos basados en atención
(Transformers) requeriría extender el módulo de generación de caminos y no puede garantizarse sin
evaluación experimental adicional.

\section{Trabajo futuro}

Las siguientes líneas de mejora se identifican como prioritarias para el trabajo futuro:

\begin{enumerate}
  \item \textbf{Selección automática del \texttt{initial\_layer}}: Diseñar un criterio cuantitativo
        (por ejemplo, basado en el número de caminos o en la complejidad del problema) que guíe la
        elección del punto de corte de forma automática, reduciendo la necesidad de conocimiento
        experto sobre la arquitectura.

  \item \textbf{Soporte para arquitecturas con conexiones residuales} (ResNet, EfficientNet): Las
        conexiones \textit{skip} rompen la estructura secuencial de caminos asumida por la
        implementación actual. Extender el módulo \texttt{paths/} para manejar estas arquitecturas
        ampliaría considerablemente el ámbito de aplicación.

  \item \textbf{Métricas de evaluación más completas}: Los experimentos actuales se centran en
        precisión global y tiempo. Añadir métricas de uso de memoria, análisis de la distribución
        de pesos de caminos, y curvas de estabilidad-plasticidad permitiría una evaluación más
        rigurosa.

  \item \textbf{Convergencia del módulo de gradiente en alta complejidad}: Explorar estrategias de
        warm-start, normalización por lotes, o ajuste adaptativo del learning rate para hacer
        converger el módulo de gradiente en escenarios de 100 o más clases.

  \item \textbf{Benchmarks adicionales}: Evaluar Qsimov sobre benchmarks estándar de aprendizaje
        continual como Split-CIFAR-100, Permuted-MNIST, o CORe50, para comparar cuantitativamente
        con métodos como EWC o iCaRL en condiciones homologables.
\end{enumerate}
